% Document/LinearAlgebra/VectorSpaces/Matrices/Part1_Matrices.tex
% Reordered Chapter 3 -- Part 0 (Why matrices?) + Part 1 (the bijection
% Hom_K(K^{(X)}, K^{(Y)}) = column-finite matrices, transport of ALL structure, examples).
% General cardinalities; column-finite matrices; sum convention \sum_{j \in X}. DRAFT for review.

\newpage
\section{Matrices}
\label{sec:matrices}

\subsection{Why Matrices?}

\begin{intuition}
    The previous chapter built linear maps as \emph{abstract} objects: a linear map $T : U \to V$ is a
    function respecting addition and scalar multiplication, and we proved theorems about such functions
    --- about kernels, images, rank, injectivity --- without ever writing a single number. That
    abstraction is exactly what makes the theory clean, but it is useless for \emph{computation}. To
    actually evaluate $T$ on a vector, to compose two maps, to invert one, or to solve $T(v) = w$, we
    need a concrete, mechanical surrogate for $T$ --- an array of scalars we can add and multiply by
    rote.

    A \textbf{matrix} is that surrogate. The whole chapter is the disciplined construction of the
    dictionary between maps and their numerical surrogates: we show (i) that for maps between
    coordinate spaces the surrogate loses \emph{no} information --- the correspondence is a bijection
    --- and (ii) that every operation on maps (composition, addition, scaling, later inversion) can be
    \emph{transported} across the dictionary into a purely arithmetic operation on the surrogates. Once
    that is done, ``computing with a linear map'' becomes ``computing with an array of scalars,'' and
    the abstract guarantees of the previous chapter certify that the computation means what we want.
\end{intuition}

\begin{goal}
    Fix the field $K$. For index sets $X$ (the \emph{source} index) and $Y$ (the \emph{target} index)
    of \emph{arbitrary} cardinality, we will:
    \begin{enumerate}
        \item define the set $\MatSpace{Y}{X}[K]$ of $(Y \times X)$-\emph{grids} of scalars --- a pure
              data structure, with no operations assumed --- and its column-finite subset
              $\CFMat{Y}{X}[K]$;
        \item exhibit a canonical \textbf{bijection} between linear maps $K^{(X)} \to K^{(Y)}$ and the
              column-finite grids $\CFMat{Y}{X}[K]$;
        \item \emph{transport} every operation --- addition, scalar multiplication, and composition ---
              from maps onto grids, and \emph{derive} the familiar entry formulas as theorems rather
              than positing them.
    \end{enumerate}
    The finite case --- $X = n = \{0, \dots, n-1\}$ for the source and $Y = m = \{0, \dots, m-1\}$ for
    the target, giving maps $K^n \to K^m$ --- is the one to keep in mind; there $K^{(X)} = K^n$ and
    every grid is automatically column-finite.
\end{goal}

\subsection{Linear Maps Between Coordinate Spaces and Their Grids}
\label{subsec:bijection-hom-mat}

Recall that for an index set $X$ the coordinate space $K^{(X)}$ is the space of \emph{finitely
supported} families $X \to K$, with canonical basis $(\CanonicalVector{j})_{j \in X}$, where
$\CanonicalVector{j}(i) = \delta_{ij}$. We begin with the data structure, deliberately stripped of all
algebraic structure.

\begin{definition}[Grids and column-finite grids]
\label{def:matrix-grid}
    Let $X, Y$ be index sets. A \textbf{$(Y \times X)$-grid over $K$} is a function
    \[
        A : Y \times X \longrightarrow K .
    \]
    We write $A_{ij} \defeq A(i, j)$ for its \textbf{entries}, where $i \in Y$ is the \textbf{row index}
    and $j \in X$ the \textbf{column index}, and $A(-, j) : i \mapsto A_{ij}$ for the $j$-th
    \textbf{column}. The set of all grids is $\MatSpace{Y}{X}[K] \defeq K^{Y \times X}$. A grid $A$ is
    \textbf{column-finite} if every column has finite support, $A(-, j) \in K^{(Y)}$ for all $j \in X$;
    the set of column-finite grids is
    \[
        \CFMat{Y}{X}[K] \defeq \Set{ A \in \MatSpace{Y}{X}[K] \mid A(-, j) \in K^{(Y)} \ \text{for all}\ j \in X } .
    \]
    At this stage these are merely \emph{sets of grids}: no addition, no multiplication, no
    interpretation as a map.
\end{definition}

\begin{remark}[{The convention $\MatSpace{Y}{X}[K]$ for maps $K^{(X)} \to K^{(Y)}$, and the finite case}]
\label{rem:convention-and-finite}
    The row index $Y$ is written first and matches the \emph{target} $K^{(Y)}$; the column index $X$
    matches the \emph{source} $K^{(X)}$. This is the convention under which composition of maps reads as
    left-to-right multiplication of grids (Theorem~\ref{thm:entry-formulas}). When $X$ and $Y$ are
    \emph{finite}, say $X = n = \{0, \dots, n-1\}$ and $Y = m = \{0, \dots, m-1\}$, every column is
    automatically finitely supported, so $\CFMat{m}{n}[K] = \MatSpace{m}{n}[K]$ is the set of ordinary
    $m \times n$ matrices and $K^{(n)} = K^n$. Column-finiteness only bites when $Y$ is infinite.
\end{remark}

\begin{remark}[Column-finiteness is weaker than global finiteness]
\label{rem:currying-cfmat}
    Reading a grid one column at a time identifies $\CFMat{Y}{X}[K]$ with the functions
    $X \to K^{(Y)}$, $j \mapsto A(-, j)$ --- one finitely supported column per index $j$ --- so that
    $\CFMat{Y}{X}[K] \cong \bigl(K^{(Y)}\bigr)^{X}$. This is \emph{not} the space $K^{(Y \times X)}$ of
    grids with \textbf{globally} finite support: column-finiteness constrains each column on its own and
    still permits infinitely many nonzero columns, whereas $K^{(Y \times X)}$ allows only finitely many
    nonzero entries in total. The two notions coincide when $X$ is finite and diverge strictly once $X$
    is infinite. It is column-finiteness --- not global finiteness --- that matches the linear maps
    $K^{(X)} \to K^{(Y)}$, as the next theorem shows.
\end{remark}

The link between a map and a grid rests on one elementary fact: a linear map out of $K^{(X)}$ is
completely pinned down by what it does to the canonical basis, and each such value lies in $K^{(Y)}$,
hence has finitely many nonzero coordinates.

\begin{theorem}[The Map--Grid Bijection]
\label{thm:map-grid-bijection}
    Let $X, Y$ be index sets. The assignment
    \[
        \Phi : \Hom[K]{K^{(X)}}{K^{(Y)}} \longrightarrow \MatSpace{Y}{X}[K],
        \qquad
        \Phi(T)_{ij} \defeq \bigl( T(\CanonicalVector{j}) \bigr)_i ,
    \]
    whose $j$-th column is the image vector $T(\CanonicalVector{j}) \in K^{(Y)}$, takes values in
    $\CFMat{Y}{X}[K]$ and is a \emph{bijection}
    \[
        \Phi : \Hom[K]{K^{(X)}}{K^{(Y)}} \xrightarrow{\ \sim\ } \CFMat{Y}{X}[K] ,
    \]
    with inverse denoted $\Phi^{-1}$.
\end{theorem}

\begin{proof}
    \emph{$\Phi$ lands in $\CFMat{Y}{X}[K]$.} For each $T$, the $j$-th column of $\Phi(T)$ is
    $T(\CanonicalVector{j}) \in K^{(Y)}$, which has finite support; hence $\Phi(T)$ is column-finite.

    \emph{Injective.} Suppose $\Phi(S) = \Phi(T)$. Then $S(\CanonicalVector{j}) = T(\CanonicalVector{j})$
    for every $j \in X$, since two vectors of $K^{(Y)}$ with equal coordinates are equal. A linear map
    is determined by its values on a basis: for arbitrary $x = \sum_{j \in X} x_j \CanonicalVector{j}
    \in K^{(X)}$ --- a finite sum, as $x$ has finite support ---
    \[
        S(x) = \sum_{j \in X} x_j\, S(\CanonicalVector{j}) = \sum_{j \in X} x_j\, T(\CanonicalVector{j}) = T(x),
    \]
    so $S = T$.

    \emph{Surjective.} Let $A \in \CFMat{Y}{X}[K]$. Each column $A(-, j)$ lies in $K^{(Y)}$, so there is
    a \emph{unique} linear map $T_A : K^{(X)} \to K^{(Y)}$ with $T_A(\CanonicalVector{j}) \defeq A(-, j)$
    (a basis determines a linear map uniquely by prescribing its values on the basis). By construction
    $\Phi(T_A)_{ij} = \bigl( T_A(\CanonicalVector{j}) \bigr)_i = A_{ij}$, i.e.\ $\Phi(T_A) = A$. Thus
    $\Phi$ is onto, and being injective and surjective it is a bijection.
\end{proof}

\begin{remark}[The fundamental datum: rows index the target, columns the source]
\label{rem:fundamental-datum}
    The double index $(i, j)$ is the central bookkeeping device of the chapter. The \textbf{column}
    index $j$ selects a source basis vector $\CanonicalVector{j}$, and the corresponding column of
    $\Phi(T)$ \emph{is} its image $T(\CanonicalVector{j})$; the \textbf{row} index $i$ then reads off
    the $i$-th coordinate of that image in the target. ``Apply $T$ to each input basis vector and stack
    the results as columns'' is the entire recipe for $\Phi(T)$.
\end{remark}

\begin{remark}[{Convention: why $(T(\CanonicalVector{j}))_i$ and not $(T(\CanonicalVector{i}))_j$?}]
\label{rem:index-swap-convention}
    The definition $\Phi(T)_{ij} = \bigl( T(\CanonicalVector{j}) \bigr)_i$ is a \emph{choice}. We could
    equally have declared the $(i, j)$-entry to be $\bigl( T(\CanonicalVector{i}) \bigr)_j$, which would
    identify $\Hom[K]{K^{(X)}}{K^{(Y)}}$ with $\MatSpace{X}{Y}[K]$ (an $X \times Y$ matrix) rather than
    with $\MatSpace{Y}{X}[K]$ (a $Y \times X$ matrix). In some sense that alternative is more natural:
    the size of the matrix would then match the order $(X, Y)$ in which source and target appear in
    $K^{(X)} \to K^{(Y)}$.

    We adopt the swapped convention $\Phi(T)_{ij} = \bigl( T(\CanonicalVector{j}) \bigr)_i$ instead,
    because we want a map $K^{(X)} \to K^{(Y)}$ to be recorded as a $Y \times X$ matrix --- the
    \emph{target} index $i$ labelling rows, the \emph{source} index $j$ labelling columns. This is the
    standard convention. Under it the image $T(\CanonicalVector{j})$ of the $j$-th source basis vector
    is exactly the $j$-th \emph{column} of $\Phi(T)$, matching the recipe of
    Remark~\ref{rem:fundamental-datum}; the alternative would stack those images as rows instead. We
    follow the target-first ($Y \times X$) convention throughout.
\end{remark}

\begin{definition}[The identity grid]
\label{def:identity-grid}
    The \textbf{identity grid} on an index set $X$ is the image of the identity map under the bijection,
    \[
        \IdentityMat{X} \defeq \Phi\bigl(\Identity{K^{(X)}}\bigr) \in \CFMat{X}{X}[K] .
    \]
    Its entries are the Kronecker delta, read straight off the definition of $\Phi$:
    \[
        (\IdentityMat{X})_{ij}
        = \bigl(\Identity{K^{(X)}}(\CanonicalVector{j})\bigr)_i
        = (\CanonicalVector{j})_i = \delta_{ij}
        = \begin{cases} 1 & i = j, \\ 0 & i \neq j, \end{cases}
    \]
    so its $j$-th column is $\CanonicalVector{j}$. When $X = n$ is finite we write $\IdentityMat{n}$ for
    the $n \times n$ identity. Once composition is transported (Subsection~\ref{subsec:transport}),
    $\IdentityMat{X}$ is the two-sided multiplicative unit (Corollary~\ref{cor:inherited-laws}) --- as
    befits the grid of the identity map.
\end{definition}

\begin{remark}[{A one-column grid is a vector: the case $X = 1$}]
\label{rem:column-grid-is-vector}
    The most degenerate source index makes the bijection reproduce a fact from the previous chapter. Take
    $X = 1 = \{0\}$, a one-element set, so that $K^{(X)} = K^{(\{0\})} = K$ with canonical basis the
    single vector $\CanonicalVector{0} = 1$. A grid in $\CFMat{Y}{1}[K]$ then has a \emph{single} column,
    and Theorem~\ref{thm:map-grid-bijection} specialises to
    \[
        \Hom[K]{K}{K^{(Y)}} \xrightarrow{\ \sim\ } \CFMat{Y}{1}[K],
        \qquad T \longmapsto T(1),
    \]
    that one column being the image $T(1) \in K^{(Y)}$. Composing with the natural isomorphism
    $\Hom[K]{K}{K^{(Y)}} \cong K^{(Y)}$ that evaluates a map at $1$
    (Theorem~\ref{thm:hom-k-v-iso}, applied to $V = K^{(Y)}$) collapses both sides to
    \[
        \CFMat{Y}{1}[K] \cong K^{(Y)} :
    \]
    a one-column grid \emph{is} a vector of $K^{(Y)}$, namely its single column, and conversely every
    vector is the one-column grid of its own coordinates. This is why a vector is recorded as a column.
    The same identification, carried out for an \emph{abstract} space $V$ through a chosen basis, is what
    makes the coordinate vector $\CoordVec[\mathcal B]{v}$ itself a column matrix.
\end{remark}

\begin{remark}[Two structures on one set: grid versus map --- insist on this]
\label{rem:grid-vs-map}
    Theorem~\ref{thm:map-grid-bijection} identifies two \emph{a priori} different kinds of object on
    the same underlying data. It is worth dwelling on the distinction, because the rest of the chapter
    consists of moving back and forth across it.
    \begin{itemize}
        \item \textbf{The map view} ($T : K^{(X)} \to K^{(Y)}$). Here the object is a \emph{morphism}:
              it can be evaluated on vectors, \emph{composed} with other maps, and it has structural
              invariants (kernel, image, rank). These are the operations of linear algebra proper ---
              they see the object as an arrow in the category of $K$-vector spaces.
        \item \textbf{The grid view} ($A \in \CFMat{Y}{X}[K]$). Here the object is a \emph{table of
              numbers}: a function on a rectangle of indices. It supports entrywise operations that
              ignore the linear-algebraic meaning entirely --- reading off a single entry, extracting a
              row or column or diagonal, the entrywise (Hadamard) product
              $(\Hadamard{A}{B})_{ij} = A_{ij} B_{ij}$, the index swap $(i, j) \mapsto (j, i)$, sparsity
              patterns, reshaping.
    \end{itemize}
    The point of $\Phi$ is that these are two \emph{coordinatisations} of one and the same information:
    the abstract map, its presentation as data, and the bridge between them. The same three-layer
    pattern governed vectors in the previous chapters --- a geometric vector $v$, its coordinate tuple
    $\CoordVec[\mathcal B]{v} = \Chart{\mathcal B}(v)$ given by the coordinate isomorphism
    $\Chart{\mathcal B}$ of an ordered basis, and the bridge $\Chart{\mathcal B}$ itself --- and it will
    reappear in full force in a later section, when we coordinatise maps between \emph{abstract} spaces.
    For maps $K^{(X)} \to K^{(Y)}$ the middle ``coordinatisation'' layer is invisible because these
    spaces already carry canonical bases; the table below records the general shape we are heading
    toward (here $\Chart{\mathcal B} : V \to K^{(X)}$ is the coordinate isomorphism of a basis
    $\mathcal B$ of $V$).

    \begin{center}
    \begin{tabular}{l|l|l|l}
        Level & Geometric object & Coordinatisation & Data structure \\
        \hline
        Scalar & $a \in K$ & $a \in K$ & a field element \\
        Vector & $v \in V$ & $\Chart{\mathcal B}(v) \in K^{(X)}$ & tuple $(a_i)_{i \in X}$ \\
        Map & $T : U \to V$ & $\Composition{\Chart{\mathcal B_V}}{\Composition{T}{\Frame{\mathcal B_U}}}$ & grid $(A_{ij})_{(i,j) \in Y \times X}$ \\
    \end{tabular}
    \end{center}

    Until we make the identification official in a later section, we keep the two views
    \emph{notationally distinct}: $\Phi(T)$ is the grid of a map, and $\Phi^{-1}(A)$ is the map of a
    grid. We do not yet speak of ``the kernel of a matrix'' or let a bare grid act on vectors; that
    licence is granted only after the identification.
\end{remark}

\subsection{Transport of Structure}
\label{subsec:transport}

The set $\Hom[K]{K^{(X)}}{K^{(Y)}}$ of linear maps is itself a $K$-vector space under the
\textbf{pointwise} operations
\[
    (S + T)(x) \defeq S(x) + T(x),
    \qquad
    (\lambda T)(x) \defeq \lambda\, T(x),
\]
which are again linear (a linear combination of linear maps is linear, as established for linear maps),
and linear maps can be \emph{composed}. Using the bijection $\Phi$ of
Theorem~\ref{thm:map-grid-bijection} --- and its inverse $\Phi^{-1}$ --- we \textbf{transport} each of
these operations onto column-finite grids: we \emph{define} the matrix operation to be whatever $\Phi$
turns the map operation into. Nothing is posited about grids directly; every entry formula below is
then \emph{derived}.

\begin{definition}[Matrix operations, by transport along $\Phi$]
\label{def:matrix-ops-transport}
    Let $\Phi$ be the bijection of Theorem~\ref{thm:map-grid-bijection} and $\Phi^{-1}$ its inverse.
    For $A, B \in \CFMat{Y}{X}[K]$ and $\lambda \in K$ define the \textbf{sum} and \textbf{scalar
    multiple} by transporting the pointwise operations on $\Hom[K]{K^{(X)}}{K^{(Y)}}$:
    \[
        A + B \defeq \Phi\bigl(\Phi^{-1}(A) + \Phi^{-1}(B)\bigr),
        \qquad
        \lambda A \defeq \Phi\bigl(\lambda\, \Phi^{-1}(A)\bigr).
    \]
    For $A \in \CFMat{Y}{Z}[K]$ and $B \in \CFMat{Z}{X}[K]$ define the \textbf{product} by transporting
    composition,
    \[
        A B \defeq \Phi\bigl(\Composition{\Phi^{-1}(A)}{\Phi^{-1}(B)}\bigr) \in \CFMat{Y}{X}[K],
    \]
    where $\Phi^{-1}(A) : K^{(Z)} \to K^{(Y)}$ and $\Phi^{-1}(B) : K^{(X)} \to K^{(Z)}$ compose to a
    map $K^{(X)} \to K^{(Y)}$, whose grid is again column-finite by
    Theorem~\ref{thm:map-grid-bijection}.
\end{definition}

Because they are transports of operations on a vector space along a bijection, these endow
$\CFMat{Y}{X}[K]$ with a $K$-vector space structure for which $\Phi$ is, \emph{by construction}, a
linear isomorphism; and the product inherits associativity, bilinearity, and units from composition.
What remains is to compute what these transported operations \emph{are} on entries. Note in particular
that addition and scalar multiplication are obtained here \emph{by transport}, not posited entrywise
--- the entrywise formulas are a conclusion.

\begin{theorem}[Entry formulas]
\label{thm:entry-formulas}
    The transported operations of Definition~\ref{def:matrix-ops-transport} are given on entries by
    \[
        (A + B)_{ij} = A_{ij} + B_{ij},
        \qquad
        (\lambda A)_{ij} = \lambda\, A_{ij},
        \qquad
        (A B)_{ij} = \sum_{k \in Z} A_{ik}\, B_{kj} ,
    \]
    where each product sum is finite because $B(-, j) \in K^{(Z)}$ has finite support, and the product
    $AB$ is again column-finite. Moreover, for $A \in \CFMat{Y}{X}[K]$ and $x \in K^{(X)}$, the map
    $\Phi^{-1}(A)$ acts by
    \[
        \bigl( \Phi^{-1}(A)(x) \bigr)_i = \sum_{j \in X} A_{ij}\, x_j ,
    \]
    a finite sum since $x$ has finite support --- the matrix--vector product.
\end{theorem}

\begin{proof}
    \emph{Action on a vector.} Write $S = \Phi^{-1}(A)$, so $S(\CanonicalVector{j}) = A(-, j)$ has
    $i$-th coordinate $A_{ij}$. For $x = \sum_{j \in X} x_j \CanonicalVector{j}$ (finite support),
    linearity gives $S(x) = \sum_{j \in X} x_j\, S(\CanonicalVector{j})$, whose $i$-th coordinate is
    $\sum_{j \in X} x_j A_{ij} = \sum_{j \in X} A_{ij} x_j$ (a finite sum).

    \emph{Sum and scalar multiple.} Let $S = \Phi^{-1}(A)$, $T = \Phi^{-1}(B)$. By definition
    $A + B = \Phi(S + T)$, so
    \[
        (A + B)_{ij} = \bigl((S + T)(\CanonicalVector{j})\bigr)_i
                     = \bigl(S(\CanonicalVector{j})\bigr)_i + \bigl(T(\CanonicalVector{j})\bigr)_i
                     = A_{ij} + B_{ij},
    \]
    using additivity of coordinates in $K^{(Y)}$. Likewise
    $(\lambda A)_{ij} = \bigl((\lambda S)(\CanonicalVector{j})\bigr)_i
    = \lambda\, \bigl(S(\CanonicalVector{j})\bigr)_i = \lambda A_{ij}$.

    \emph{Product.} Let $S = \Phi^{-1}(A) : K^{(Z)} \to K^{(Y)}$ and $T = \Phi^{-1}(B) : K^{(X)} \to
    K^{(Z)}$. By definition $A B = \Phi(\Composition{S}{T})$, whose $j$-th column is
    $(\Composition{S}{T})(\CanonicalVector{j}) = S\bigl(T(\CanonicalVector{j})\bigr) = S\bigl(B(-,j)\bigr)$.
    Now $B(-, j) = \sum_{k \in Z} B_{kj}\, \CanonicalVector{k}$ is a finite sum (column-finiteness of
    $B$), so by linearity of $S$ and the action formula just proved,
    \[
        (A B)_{ij} = \Bigl( S\bigl(B(-, j)\bigr) \Bigr)_i
        = \sum_{k \in Z} B_{kj}\, \bigl( S(\CanonicalVector{k}) \bigr)_i
        = \sum_{k \in Z} B_{kj}\, A_{ik}
        = \sum_{k \in Z} A_{ik}\, B_{kj} .
    \]
    Finally, $\Set{ i \in Y \mid (AB)_{ij} \neq 0 } \SubsetEq \BigUnion[k \in \operatorname{supp} B(-,j)]
    \Set{ i \in Y \mid A_{ik} \neq 0 }$ is a finite union of finite sets, so $(AB)(-, j) \in K^{(Y)}$ and $AB$ is
    column-finite, consistent with Definition~\ref{def:matrix-ops-transport}.
\end{proof}

\begin{remark}[The matrix--vector product is transport of structure too]
\label{rem:matvec-is-transport}
    The action formula $\bigl(\Phi^{-1}(A)(x)\bigr)_i = \sum_{j \in X} A_{ij}\, x_j$ reads like a new
    operation --- a matrix acting on a vector --- but it is transport of structure once more, with
    nothing extra to prove. By Remark~\ref{rem:application-is-composition}, applying the map
    $\Phi^{-1}(A)$ to a vector $x \in K^{(X)}$ is composition with the name $T_x : K \to K^{(X)}$,
    $1_K \mapsto x$, of that vector: $\Phi^{-1}(A)(x) = \Composition{\Phi^{-1}(A)}{T_x}$. Read through
    $\Phi$, the vector $x$ is the single-column grid $\Phi(T_x)$ --- the one-column-grid-is-a-vector
    identification of Remark~\ref{rem:column-grid-is-vector}, the case $X = 1$ where
    $\Hom[K]{K}{K^{(X)}} \cong K^{(X)}$ --- and the matrix--vector product $A x$ is the matrix product of
    $A$ with that column (Definition~\ref{def:matrix-ops-transport}). Hence ``matrix times vector'' is not
    a separate primitive: it is matrix multiplication with a one-column factor --- a special case of the
    matrix product --- exactly as application is the one-argument case of composition.
\end{remark}

\begin{corollary}[Inherited algebraic laws]
\label{cor:inherited-laws}
    Matrix addition and scalar multiplication make $\CFMat{Y}{X}[K]$ a $K$-vector space, and $\Phi$ is
    a linear isomorphism $\Hom[K]{K^{(X)}}{K^{(Y)}} \xrightarrow{\sim} \CFMat{Y}{X}[K]$. Matrix
    multiplication is associative and bilinear wherever the index sets match, with the identity grid
    $\IdentityMat{X}$ (Definition~\ref{def:identity-grid}) as two-sided unit:
    \[
        (AB)C = A(BC), \quad A(B + C) = AB + AC, \quad (A + B)C = AC + BC,
        \quad \IdentityMat{Y} A = A = A\, \IdentityMat{X} .
    \]
    In particular, for $X = Y$, $\Phi$ is an isomorphism of $K$-algebras
    $\End[K]{K^{(X)}} \xrightarrow{\sim} \CFMat{X}{X}[K]$.
\end{corollary}

\begin{proof}
    Each law holds for the corresponding operation on maps --- the pointwise operations make
    $\Hom[K]{K^{(X)}}{K^{(Y)}}$ a vector space, composition is associative and bilinear, and the
    identity map $\Identity{K^{(X)}}$ is its unit. Transporting along the bijection $\Phi$ preserves
    every such identity: associativity follows from
    $\Composition{(\Composition{R}{S})}{T} = \Composition{R}{(\Composition{S}{T})}$ together with the
    definitions, and $\Phi(\Identity{K^{(X)}}) = \IdentityMat{X}$ since
    $\bigl(\Identity{K^{(X)}}(\CanonicalVector{j})\bigr)_i = (\CanonicalVector{j})_i = \delta_{ij}$. No
    entrywise computation is needed.
\end{proof}

\begin{remark}[The general grid space and the action map]
\label{rem:action-map}
    Dropping column-finiteness, every grid $A \in \MatSpace{Y}{X}[K]$ still defines a linear
    \textbf{action map}
    \[
        \mu_A : K^{(X)} \to K^{Y}, \qquad \bigl(\mu_A(x)\bigr)_i \defeq \sum_{j \in X} A_{ij}\, x_j ,
    \]
    the sum being finite because $x \in K^{(X)}$ has finite support. Column-finiteness of $A$ is
    exactly the condition that $\mu_A$ land in $K^{(Y)}$ rather than the full product $K^Y$, i.e.\ that
    $A$ represent a map $K^{(X)} \to K^{(Y)}$. Thus $\CFMat{Y}{X}[K]$ is the correct home for matrices
    of maps between coordinate spaces at \emph{any} cardinality; the larger $\MatSpace{Y}{X}[K]$ becomes
    indispensable only once a topology lets one make sense of non-finite sums (the passage from
    algebra to functional analysis).
\end{remark}

\subsection{Practical Examples}

We compute with the formulas just derived; everything below is elementary arithmetic in the finite case
$K^n \to K^m$, where every grid is column-finite.

\begin{example}[Reading off $\Phi(T)$]
\label{ex:read-off}
    Let $T : K^3 \to K^2$ be $T(x_0, x_1, x_2) = (x_0 + 2 x_2,\; x_0 - x_1)$. Applying $T$ to the
    canonical basis $(\CanonicalVector{0}, \CanonicalVector{1}, \CanonicalVector{2})$ of $K^3$ gives the
    columns:
    \[
        T(\CanonicalVector{0}) = (1, 1), \quad
        T(\CanonicalVector{1}) = (0, -1), \quad
        T(\CanonicalVector{2}) = (2, 0),
        \qquad
        \Phi(T) = \begin{pmatrix} 1 & 0 & 2 \\ 1 & -1 & 0 \end{pmatrix} \in \CFMat{2}{3}[K].
    \]
    The action formula recovers $T$: for $x = (x_0, x_1, x_2)$,
    $\bigl(\Phi^{-1}(\Phi(T))(x)\bigr) = \bigl(\sum_{j \in 3} \Phi(T)_{0j} x_j,\ \sum_{j \in 3}
    \Phi(T)_{1j} x_j\bigr) = (x_0 + 2x_2,\; x_0 - x_1) = T(x)$.
\end{example}

\begin{example}[Composition is multiplication]
\label{ex:composition-is-mult}
    Let $T : K^3 \to K^2$ be as above and $S : K^2 \to K^2$, $S(y_0, y_1) = (y_1,\; y_0 + y_1)$, so
    $\Phi(S) = \begin{psmallmatrix} 0 & 1 \\ 1 & 1 \end{psmallmatrix}$. The composite
    $\Composition{S}{T} : K^3 \to K^2$ has grid
    \[
        \Phi(\Composition{S}{T}) = \Phi(S)\, \Phi(T)
        = \begin{pmatrix} 0 & 1 \\ 1 & 1 \end{pmatrix}
          \begin{pmatrix} 1 & 0 & 2 \\ 1 & -1 & 0 \end{pmatrix}
        = \begin{pmatrix} 1 & -1 & 0 \\ 2 & -1 & 2 \end{pmatrix}.
    \]
    Checking directly: $S(T(x)) = S(x_0 + 2x_2,\, x_0 - x_1) = (x_0 - x_1,\; 2x_0 - x_1 + 2x_2)$, whose
    grid is exactly the product above. The shared index $k$ in
    $(\Phi(S)\Phi(T))_{ij} = \sum_{k \in 2} \Phi(S)_{ik}\Phi(T)_{kj}$ is the intermediate space $K^2$
    being cancelled --- this is why the inner index sets must agree.
\end{example}

\begin{example}[Associativity for free]
    With $A = \begin{psmallmatrix} 1 & 1 \\ 0 & 1 \end{psmallmatrix}$,
    $B = \begin{psmallmatrix} 2 & 0 \\ 1 & 3 \end{psmallmatrix}$,
    $C = \begin{psmallmatrix} 1 \\ 1 \end{psmallmatrix}$ one verifies
    $(AB)C = \begin{psmallmatrix} 3 & 3 \\ 1 & 3 \end{psmallmatrix}\begin{psmallmatrix} 1 \\ 1 \end{psmallmatrix}
    = \begin{psmallmatrix} 6 \\ 4 \end{psmallmatrix}$ and
    $A(BC) = \begin{psmallmatrix} 1 & 1 \\ 0 & 1 \end{psmallmatrix}\begin{psmallmatrix} 2 \\ 4 \end{psmallmatrix}
    = \begin{psmallmatrix} 6 \\ 4 \end{psmallmatrix}$. The agreement is guaranteed in advance by
    Corollary~\ref{cor:inherited-laws}: it is associativity of composition, read through $\Phi$.
\end{example}
